\documentclass[12pt,a4paper]{config/myConfig}
\usepackage{silence}
\WarningFilter{latex}{You have requested}
\usepackage{config/myDependencies}
\usepackage{style/myEnv}
\usepackage{style/myCommands}
\usepackage{style/myColors}
\usepackage{hyperref}
\usepackage{parskip}

\addbibresource{bibliography.bib}  % your .bib file

%Number of the practica for this assigment:
\newcommand{\AssigmentNumber}{Laboratorio 1}%<- Nombre de la Entrega
\newcommand{\AssigmentDate}{6 Febrero, 2026} %<- Fecha de Entrega
%------------------
%If any computer language used:
\newcommand{\Rcode}{R - versión 3.4.4}
\newcommand{\Pycode}{Python - versión 3.13}
\newcommand{\Trackercode}{Tracker - versión 5.0.1}

%----------------------
%Palabras Clave para el abstract del informe:
\newcommand{\PalabrasClave}{\small \it
..,$\ $.. ,$\ $ ..}

\title{\Large $\ $\\ \bf Velocidad terminal en caída libre a través de un fluido 
viscoso} %<-Titulo del informe

\author
\info

\abstract{\PalabrasClave
\normalsize
\\[17pt]
{\bf Abstract.} 
En este trabajo se estudió experimentalmente la velocidad terminal de esferas que caen en aceite vegetal, con el fin de analizar cómo se relaciona la masa del objeto con la velocidad terminal que alcanza. A diferencia de los experimentos tradicionales en aire, donde el arrastre suele modelarse como proporcional al cuadrado de la velocidad, aquí se trabajó en un fluido más viscoso para favorecer un régimen donde el arrastre es aproximadamente lineal y el comportamiento resulta más estable.

Se utilizaron tres canicas (dos de vidrio y una de metal) con volúmenes similares pero masas distintas. El movimiento se grabó en video y los datos de posición–tiempo se obtuvieron con \textit{Tracker}. A partir de estos datos se estimó la velocidad terminal mediante regresiones lineales agrupando las réplicas, y luego se evaluaron los supuestos estadísticos del modelo.

Los resultados muestran que la canica de mayor masa (16.311 g) presentó un comportamiento más cercano al modelo teórico, alcanzando con mayor claridad un régimen aproximadamente lineal en la gráfica posición–tiempo. En cambio, las canicas más livianas no llegaron completamente a la velocidad terminal dentro del intervalo observado, lo que produjo aceleración residual y fallas en los supuestos del modelo, reduciendo la confiabilidad de las estimaciones.

En general, no se pudo confirmar con suficiente solidez la relación lineal entre masa y velocidad terminal bajo las condiciones del experimento. Sin embargo, la metodología mostró ser más confiable cuando se emplean masas mayores. Este estudio propone una alternativa didáctica para analizar el arrastre viscoso en laboratorio y resalta la importancia de elegir adecuadamente las masas para obtener resultados consistentes.
}

%Table specifications:
\sisetup{
  separate-uncertainty = true,
  table-number-alignment = center,
  detect-weight = true,
  detect-family = true
}

\begin{document}
\thispagestyle{myheadings}
\pagestyle{myheadings}
\markright{\tt \AssigmentNumber|    \AssigmentDate}%check the date



\section{Marco Conceptual}
\label{sec:MARCO-CONCEPTUAL}
\subsection{Marco teórico}
\label{sec:MacroTeorico}

El arrastre es la fuerza que se opone al movimiento relativo de un objeto a través de un fluido. En general, dicha fuerza puede expresarse como proporcional a alguna potencia de la velocidad del objeto y a un coeficiente que depende de las propiedades del fluido y de la geometría del cuerpo. La forma específica que adopta el arrastre depende críticamente del número de Reynolds, definido como

\[
Re = \frac{\rho v D}{\mu},
\]

donde $\rho$ es la densidad del fluido, $v$ la velocidad característica del objeto, $D$ una longitud característica (para una esfera, su diámetro), y $\mu$ la viscosidad dinámica del fluido. Este número adimensional cuantifica la relación entre las fuerzas inerciales y las fuerzas viscosas que actúan sobre el sistema.

Cuando $Re \ll 1$, las fuerzas viscosas dominan sobre las inerciales y el flujo alrededor del objeto es altamente ordenado y laminar. En este régimen, conocido como régimen de Stokes, el arrastre es proporcional linealmente a la velocidad:

\[
F_D = 6\pi \mu r v,
\]

para una esfera de radio $r$. Bajo estas condiciones, la velocidad terminal se obtiene cuando el arrastre equilibra el peso efectivo (considerando el empuje), resultando en una dependencia lineal entre la fuerza de arrastre y la velocidad.

Por el contrario, cuando $Re \gg 1$, las fuerzas inerciales predominan y el flujo puede presentar separación y formación de estelas. En este régimen, el arrastre es proporcional al cuadrado de la velocidad:

\[
F_D = \frac{1}{2} C_D \rho A v^2,
\]

donde $C_D$ es el coeficiente de arrastre y $A$ el área transversal del objeto.

En la presente investigación, dado que se emplea un fluido con viscosidad significativamente mayor que la del aire (aceite vegetal), el número de Reynolds se reduce respecto al caso aerodinámico clásico. Bajo estas condiciones, el comportamiento del sistema se aproxima a un régimen dominado por efectos viscosos, por lo que se adopta el modelo de arrastre lineal como marco teórico principal para el análisis del movimiento. \citet{Saldania2024The}

Considérese un cuerpo de geometría esférica, con masa $m$ y volumen $V$, que cae
bajo su propio peso en un fluido con arrastre lineal. La sumatoria de fuerzas
sobre el cuerpo es entonces:
\begin{equation*}
    \sum F = mg - kv-B = ma
\end{equation*}
donde $k$ es el coeficiente de arrastre y $B=\rho g V$ es la fuerza boyante,
con $\rho$ la densidad del fluido. Resolviendo la ecuación diferencial, se tiene que la velocidad está dada por:
\begin{equation}
    v(t) = \frac{mg-B}{k}+\left[ \frac{v_0k-gm-B}{k} \right]e^{-\frac{k}{m}t} \tag{Velocidad en función del tiempo}
\end{equation}
Notar que la función de velocidad tiene una asíntota, es decir, existe una velocidad terminal: 
\begin{equation}
    v(t\rightarrow \infty) = v_t=\frac{mg-B}{k} \tag{Velocidad terminal} 
\end{equation}
Integrando la función de velocidad se obtiene la de posición: 
\begin{equation}
    x(t) = v_tt + (v_t-v_0)\frac{m}{k}e^{-\frac{k}{m}t} \tag{Posición en función del tiempo} \label{eq:x-t}
\end{equation}

%HACER LA CORRECCIÓN EN LA DEDUCCIÓN.
Esto contrasta con la velocidad terminal para un arrastre cuadrático, cuya solución analítica es:
\begin{equation}
    v(t) = \sqrt{\frac{F}{k}}\tanh { \sqrt{\frac{Fk}{m^2}t} } \tag{Velocidad en función del tiempo}
\end{equation}
donde $F$ es la fuerza de empuje, igual al peso en el aire y 
al peso menos la boyante en fluidos más densos. En este caso, la velocidad terminal es: 
\begin{equation}
    v(t\rightarrow \infty) = v_t=\sqrt{\frac{F}{k}} \tag{Velocidad terminal} 
\end{equation}
Integrando la función de velocidad se obtiene la de posición: 
\begin{equation}
    x(t) = \frac{m}{k} \sqrt{v_t} \ln \left[ \cosh(\frac{k}{m}\sqrt{v_t} t) \right] \tag{Posición en función del tiempo}
\end{equation}

\subsection{Antecedentes}
\label{sec:ANTECEDENTES}

\indent En \textit{Physics Teacher} se presenta un experimento propuesto por \citet{gluck2003air}, que expone una experimentación para determinar el coeficiente de resistencia en caída libre usando una pelota inflada y un globo esférico inflado. \citeauthor{gluck2003air} sugiere comprobar si hay linealidad entre la velocidad terminal y la raíz de la masa del objeto, ya que en caída libre en aire se utiliza un modelo de arrastre cuadrático. 

\indent La metodología consiste en grabar cinco globos con diferentes masas adjuntadas en caída libre, y tomar una muestra de cinco repeticiones para cada globo. De las grabaciones se obtienen los datos de altura y tiempo. \citeauthor{gluck2003air} sugiere implementar una diferenciación numérica para estimar la velocidad a partir de la relación altura-tiempo. Luego, considerando únicamente los segmentos de la gráfica velocidad-tiempo que sean aproximadamente constantes, se aplica un ajuste por mínimos cuadrados a los datos de posición tiempo en estos segmentos de tiempo. El coeficiente de esta regresión, que corresponde a una estimación de la velocidad terminal, permite finalmente obtener la constante de resistencia. 

\indent Se identificaron vulnerabilidades de errores sistemáticos y aleatorios que pueden surgir en \citet{gluck2003air}, como que la obtención de la velocidad mediante diferenciación numérica de los datos de posición amplifica el ruido instrumental del sensor ultrasónico, afectando la estabilidad de la región identificada como velocidad terminal. Asimismo, la selección visual de los segmentos “constantes” para realizar el ajuste lineal puede introducir sesgos subjetivos en la estimación de la pendiente. A ello se suman posibles efectos no controlados, tales como corrientes de aire, derivas laterales del globo, pequeñas imprecisiones geométricas en la medición de la altura y limitaciones en la tasa de muestreo del instrumento, factores que pueden alterar tanto la determinación de la velocidad terminal como la verificación experimental del modelo $R = b v^2$.


\indent Una propuesta alternativa es utilizar canicas en aceite comercial, en lugar de globos en aire. Esta se fundamenta en consideraciones hidrodinámicas relacionadas con el número de Reynolds y la estabilidad del flujo. \citet{Xia2021} demostró que el régimen de caída de una gota en otro fluido depende críticamente del número de Reynolds: para valores elevados aparecen vórtices y separación de flujo, lo que modifica significativamente la fuerza de arrastre. En contraste, cuando el número de Reynolds disminuye, el flujo se vuelve más ordenado y los efectos viscosos adquieren mayor relevancia, haciendo que el comportamiento del arrastre sea más predecible. Dado que la viscosidad del aceite comercial es varios órdenes de magnitud mayor que la del aire, el movimiento de una esfera sólida en dicho medio ocurre en un régimen con mayor influencia viscosa que el caso aerodinámico estudiado por Gluck, reduciendo la formación de estructuras vorticiales y favoreciendo una mejor concordancia entre teoría y experimento. Este criterio resulta consistente con los resultados reportados por \cite{Ruan2021}, quienes muestran que los modelos hidrodinámicos presentan mejor ajuste cuando las fuerzas viscosas tienen un papel dominante. % add article supporting the use of liquids and provide evidence on why this is better. (DONE EHVM)
\\

Si bien \citeauthor{gluck2003air} propone un marco pedagógico sólido para el estudio de la resistencia del aire mediante objetos esféricos ligeros, su planteamiento se desarrolla principalmente en un medio gaseoso de baja viscosidad, donde los efectos inerciales y las perturbaciones ambientales pueden introducir mayores incertidumbres experimentales. El presente trabajo no busca reemplazar dicha propuesta, sino complementarla mediante el estudio de un fenómeno análogo en un medio líquido de mayor viscosidad. Al emplear esferas sólidas en aceite comercial, el sistema se traslada hacia un régimen con mayor influencia viscosa en comparación con el caso aerodinámico, reduciendo la importancia relativa de los efectos inerciales y favoreciendo un comportamiento más estable del flujo. Bajo estas condiciones, el análisis del arrastre puede realizarse con mayor control experimental y reproducibilidad. Dicha modificación vincula el enfoque didáctico de Gluck con la literatura contemporánea sobre arrastre viscoso (por ejemplo, \citeauthor{Xia2021}; \citeauthor{Ruan2021}), extendiendo el marco experimental hacia un régimen más directamente comparable con modelos hidrodinámicos bien establecidos. En este sentido, el presente estudio aporta a la literatura una alternativa experimental controlada para el análisis de las leyes de arrastre, preservando al mismo tiempo el espíritu formativo de la propuesta original.  % A final observation on what would this approach add to the literature (and which literature specficialy), but always show some degree of support to Gluck work and how his proposal will be useful with this addition to the literature. 



\section{Introducción}
\label{sec:INTRODUCCION}
\subsection{Planteamiento del Problema}
El problema considerado es el de una masa esférica que cae en un fluido viscoso,
específicamente aceite vegetal, bajo la acción de su peso. Además, se consideran 
una fuerza de arrastre proporcional a la velocidad, y la fuerza 
boyante. Se busca investigar la relación entre la masa en caída libre bajo estas
condiciones y la velocidad terminal que alcanza, estimando la velocidad terminal mediante una regresión lineal de los datos de posición-tiempo.
\label{sec:PLANTAMIENTO-PROBLEMA}

\subsection{Supuestos de la Investigación}
\label{sec:SUPUESTOS}

\begin{supuestos}
  \item \label{su:uno} La turbulencia en el flujo del aceite no produce fuerzas 
  significativas sobre la canica. 
  \item \label{su:dos} La velocidad inicial al dejar caer cada canica es 0.
  \item \label{su:tres} El volumen de las 3 canicas es lo suficientemente similar para asegurar que las diferencias en la velocidad terminal se deben únicamente 
  a la masa. 
\end{supuestos}

\section{Hipótesis y Objetivos}
\label{sec:HIPOTESIS-OBJECTIVOS}

\subsection{Hipótesis}
\label{sec:HIPOTESIS}

%Constante de amortiguamiento
\begin{hipotesis}
  {En una caída libre en aceite vegetal, hay una relación lineal entre la velocidad terminal $v_{\text{t}}$ y la masa $m$ del objeto con una significancia del $5$\%. } %<- La hiptesis escrita en palabras.
  {$\beta_1 = 0$} %<- hipotesis nula
  {$\beta_1 \neq 0$} %<- hipotesis alternativa
  {$v_{\text{t},i} = \beta_0 + \beta_1(m_i) + \epsilon_i$ } %<- Modelo plantiado para probar la hipotesis
  {$\epsilon_1 \cdots \epsilon_n \sim \mathcal{N}(0,\sigma^2)$} %<- Suposuciones para el Modelo como tal.
\end{hipotesis}


\subsection{Objetivos}
\label{sec:OBJECTIVOS}

\begin{objectives}
  \item \label{obj:A} Obtener experimentalmente la velocidad terminal de cada una de las canicas en caída libre en aceite vegetal mediante una regresión lineal.
  \item \label{obj:B} Aceptar o rechazar la hipótesis propuesta.
  \item \label{obj:C} Validar la metodología propuesta para la determinación de velocidades terminales en caída libre..
\end{objectives}

\section{Metodología}
\label{sec:METODOLOGIA}



\subsection{Metodos}
\label{sec:METODOS}

\indent El montaje experimental se realizó utilizando el artefacto mostrado en la \cref{fig:Montaje 1}, el cual consiste en un difusor led (tubo transparente) y una base estructural de madera. Dicho dispositivo, originalmente diseñado para la medición de densidad de fluidos, fue adaptado para el estudio del movimiento de esferas en un medio viscoso. La estructura fue donada, y presenta un sistema de soporte rígido con patas estabilizadoras de 21.5 cm de longitud. 


\indent El difusor fue dispuesto en posición completamente vertical y fijado firmemente a la base de madera, asegurando su estabilidad mecánica durante todo el procedimiento experimental. Se verifico que el eje longitudinal del tubo fuera perpendicular al plano horizontal, con el objetivo de minimizar desviaciones laterales en la trayectoria de las esferas. Previo al inicio de las mediciones, el tubo fue llenado con aceite vegetal comercial, utilizado como fluido viscoso para el experimento. El aceite fue introducido cuidadosamente hasta alcanzar un nivel adecuado que permitiera el recorrido completo de las esferas sin que estas impactaran directamente contra el fondo del tubo durante los primeros instantes de aceleración. 


\indent Se emplearon 2 canicas de vidrio y una de metal, las cuales fueron utilizados como cuerpos de prueba. Cada esfera fue depositada manualmente en reposo justo en la superficie, procurando que la velocidad inicial fuera aproximadamente nula. Tal procedimiento buscó reducir la influencia de impulsos adicionales o componentes horizontales en la trayectoria. Además, la liberación se realizo con el cuidado necesario para evitar perturbaciones innecesarias en el fluido. 


\indent El registro del movimiento se efectúo mediante un teléfono celular \textit{Samsung A22}, que destaca por su sensor principal de $48$ MP con Estabilización Óptica de Imagen (OIS), montado sobre un trípode. La cámara fue colocada a una distancia aproximada de $1.5 m$ del montaje, alineada lateralmente en un plano paralelo al plano de caída de las esferas. Asimismo, se procuró que el eje óptico de la cámara coincidiera aproximadamente con la mitad de la altura del tubo, con el fin de reducir errores de perspectiva y paralaje en el análisis de video. 


\indent Finalmente, se verificó que el encuadre incluyera la totalidad de la trayectoria vertical de las esferas dentro del tubo antes de iniciar cada grabación. Con el sistema correctamente alineado, el fluido estabilizado y la calibración definida, se dio inicio a la fase de experimentación descrita en el procedimiento general de la \cref{fig:Metodología}. Se realizaron 3 mediciones para cada canica.

\begin{FigureLab}{Diagrama de Experimentación}{}{fig:Metodología}
      \includegraphics[width=1\linewidth]{images/practica_1/diagrama_lab_1.png}
\end{FigureLab}


\subsection{Presentación de Montaje}
\label{sec:MONTAJE}

\begin{FigureLab}{Experimento Montado 1}{}{fig:Montaje 1}
      \includegraphics[width=0.19\linewidth]{images/practica_1/Experimento_Montado_2.png}
\end{FigureLab}


%En el texto: Como se muestra en la \cref{fig:Montaje 1}, \FigRef{fig:Montaje 1} 


\subsection{Materiales}
\label{sec:MATERIALES}

\begin{TableLab}{Tabla de Materiales}{}{tab:Materiales}
    \rowcolors{2}{gray!10}{white}
    \newcolumntype{Y}{>{\raggedright\arraybackslash}X}
    \begin{tabularx}{\textwidth}{@{} l l >{\hsize=0.6\hsize}Y  >{\raggedright\arraybackslash}X @{}}
        \toprule
        \rowcolor{lightOrange}
        \textbf{Material} & \textbf{Cantidad} & \textbf{Especificación} \\
        \midrule
        Aceite Vegetal & 1.6 $L$ & Detalles del aceite en el anexo (\cref{figA:tabla aceite}).  \\
        Base de madera & 1 & Uso para sostener el tubo plástico. Donde las estabilizadores miden 21.5 cm.  \\
        Canica blanca & 1 & Con diámetro: 14.12 $mm$\\
        Canica de metal  & 1 & Con diámetro: 13.92 $mm$   \\
        Canica negra & 1 & Con diámetro: 13.92 $mm$\\
        %Canica verde & 1 & Con diámetro: 14.30 $mm$\\
        Teléfono & N/a & Uso para grabar. Videos fueron grabados en 60 fps.\\
        Tracker & N/a & Uso para obtener los datos de Distancia y Altura. $\quad \quad$ Versión  5.0.1\\
        Tubo plástico & 1 & Una altura de 1.15 metros. Circunferencia de 2.5 $cm$.  \\
        \bottomrule
    \end{tabularx}
\end{TableLab}

%\cref{tab:Materiales} or \TabRef{tab:Materiales}



\section{Ecuaciones}
\label{sec:ECUACIONES}


%Example of an equation
\begin{EquationLab}
{ v_t=\frac{mg-\rho g V}{k} } % <- ecuacion con simbologia matematica
{
\begin{itemize}[label=$\circ$, left=0pt]
    \item $v_t$ es la velocidad terminal
    \item $m$ es la masa
    \item $g$ es la aceleración gravitatoria
    \item $\rho$ es la densidad del fluido
    \item $V$ es el volumen
    \item $k$ es el coeficiente de arrastre
    
\end{itemize}
} % <- Indicacion de cada simbolo puesto en la ecuacion.
\end{EquationLab}

%\cref{eq:1}

\section{Datos}
\label{sec:DATOS}

%The physical parameters or Constants that where used for the experiment:
\begin{itemize}[label=$\rightarrow$, left=0pt]
 \item Densidad del aceite : $0.9$ $\left( \dfrac{g} {cm^{3}} \right)$\hfill(Laboratorio)
 \item Masa de canica blanca : $5.316$  $\left( g \right)$\hfill(Laboratorio)
 \item Masa de canica metal : $16.311$  $\left( g \right)$\hfill(Laboratorio)
 \item Masa de canica negra : $5.324$  $\left( g \right)$\hfill(Laboratorio)
 %\item Masa de canica verde : $5.538$  $\left( g \right)$\hfill(Laboratorio)
  
  %\item para1 : num $\left( unit \right)$\hfill(\cite{})
  %Parameter, Value, Unit, Reference where obtained.
\end{itemize}


\section{Resultados}
\label{sec:RESULTADOS}


\subsection*{Resultados de Estimación}


\begin{TableLab}{Resultado de Regresión Lineal para cada Canica}{\Rcode}{}
\begin{tabular}{lccc}
\hline\hline
\multicolumn{4}{c}{\textit{Regresión Lineal con Efectos Fijos por Grupo: } $y_{it} = \alpha_0 + v_t t_{ij} + \varepsilon_{ij}$}\\
\hline\hline
\multicolumn{4}{c}{\textit{Variable dependiente: Altura ($y$ m)}}\\
\hline
 & Blanca & Negra & Metal \\
\hline
Tiempo ($t$ s) 
  & -0.1963$^{***}$ & -0.1634$^{***}$ & -0.6049$^{***}$ \\
  & (0.0002) & (0.0005) & (0.0004) \\

Segunda Replica 
  & 0.0055$^{***}$ & 0.0059$^{***}$ & -0.0015$^{***}$ \\
  & (0.0007) & (0.0019) & (0.0005) \\

Tercera Replica 
  & -0.0239$^{***}$ & -0.0448$^{***}$ & -0.0069$^{***}$ \\
  & (0.0007) & (0.0019) & (0.0005) \\

Constante  
  & -0.0230$^{***}$ & 0.9742$^{***}$ & 0.0061$^{***}$ \\
  & (0.0007) & (0.0019) & (0.0005) \\

\hline
Observaciones 
  & 540 & 505 & 162 \\

R$^{2}$ 
  & 0.9996 & 0.9959 & 0.9999 \\

F Statistic 
  & 436,003.10$^{***}$ (3;536) 
  & 40,498.24$^{***}$ (3;501) 
  & 692,889.80$^{***}$ (3;158) \\
\hline\hline
\multicolumn{4}{r}{El numero abajo del coeficiente es el error estándar.  \textit{Note:} $^{*}p<0.1$; $^{**}p<0.05$; $^{***}p<0.01$}
\end{tabular}    
\begin{tabular}[t]{lllrrc}
\toprule
Model (Blanca) & Assumption & Diagnostic & Statistic & p\_value & Verdict\\
\midrule
\cellcolor{gray!10}{Blanca (Pooled)} & \cellcolor{gray!10}{Normality} & \cellcolor{gray!10}{Shapiro–Wilk W} & \cellcolor{gray!10}{0.9569} & \cellcolor{gray!10}{0} & \cellcolor{gray!10}{Violation}\\
Blanca (Pooled) & Homoscedasticity & Breusch–Pagan X² & 71.1600 & 0 & Violation\\
\cellcolor{gray!10}{Blanca (Pooled)} & \cellcolor{gray!10}{No Autocorrelation} & \cellcolor{gray!10}{Durbin–Watson D} & \cellcolor{gray!10}{0.1052} & \cellcolor{gray!10}{0} & \cellcolor{gray!10}{Violation}\\
\bottomrule
\end{tabular}

\begin{tabular}[t]{lllrrc}
\toprule
Model (Negra) & Assumption & Diagnostic & Statistic & p\_value & Verdict\\
\midrule
\cellcolor{gray!10}{Negra (Pooled)} & \cellcolor{gray!10}{Normality} & \cellcolor{gray!10}{Shapiro–Wilk W} & \cellcolor{gray!10}{0.9816} & \cellcolor{gray!10}{0} & \cellcolor{gray!10}{Violation}\\
Negra (Pooled) & Homoscedasticity & Breusch–Pagan X² & 100.2601 & 0 & Violation\\
\cellcolor{gray!10}{Negra (Pooled)} & \cellcolor{gray!10}{No Autocorrelation} & \cellcolor{gray!10}{Durbin–Watson D} & \cellcolor{gray!10}{0.0476} & \cellcolor{gray!10}{0} & \cellcolor{gray!10}{Violation}\\
\bottomrule
\end{tabular}
\begin{tabular}[t]{lllrrc}
\toprule
Model (Metal)& Assumption & Diagnostic & Statistic & p\_value & Verdict\\
\midrule
\cellcolor{gray!10}{Metal (Pooled)} & \cellcolor{gray!10}{Normality} & \cellcolor{gray!10}{Shapiro–Wilk W} & \cellcolor{gray!10}{0.9014} & \cellcolor{gray!10}{0.0000} & \cellcolor{gray!10}{Violation}\\
Metal (Pooled) & Homoscedasticity & Breusch–Pagan X² & 11.5311 & 0.0092 & Violation\\
\cellcolor{gray!10}{Metal (Pooled)} & \cellcolor{gray!10}{No Autocorrelation} & \cellcolor{gray!10}{Durbin–Watson D} & \cellcolor{gray!10}{1.2441} & \cellcolor{gray!10}{0.0000} & \cellcolor{gray!10}{Violation}\\
\bottomrule
\end{tabular}
\end{TableLab}

\subsection*{Gráficos de Diagnósticos}
\label{sec:Graficos}

\begin{FigureLab}{Diagnósticos del Modelo Lineal para Canica de Metal}{\Rcode}{} 
\includegraphics[width=0.8\linewidth]{images/practica_1/Plot_metal.png}    
\end{FigureLab}

\begin{FigureLab}{Diagnósticos del Modelo lineal para Canica Negra}{\Rcode}{}
\includegraphics[width=0.8\linewidth]{images/practica_1/plots_negra.png}    
\end{FigureLab}

\begin{FigureLab}{Diagnósticos del Modelo lineal para Canica Blanca}{\Rcode}{}
\includegraphics[width=0.8\linewidth]{images/practica_1/plots_blanca.png}        
\end{FigureLab}




\section{Discusión de Resultados}
\label{sec:DISCUSSION}
    
Como se planteo en el \secref{sec:MARCO-CONCEPTUAL}, este experimento es complementario al estudio de \citet{gluck2003air} para estudiar la relación entre la velocidad terminal y la masa de un objecto en caída libre como fue propuesta en \secref{sec:HIPOTESIS}. Hay que resaltar nuevamente que la metodología propuesta en \secref{sec:METODOLOGIA} es distinta a la de \citeauthor{gluck2003air}. Por lo que nuestros resultados no se pueden usar para garantizar la consistencia de la experimentación por \citeauthor{gluck2003air}; pero si son complementarios en la literatura presentada en \textit{Physics Teacher} de experimentos para estudiar el régimen laminar para estudiantes de pre-grado.\newline

\indent Siguiendo la \secref{sec:INTRODUCCION}, el problema central de este estudio consiste en determinar experimentalmente la velocidad terminal de esferas sólidas en caída vertical dentro de un medio viscoso, y examinar la validez de los supuestos físicos que describen dicho fenómeno. Bajo el marco conceptual desarrollado previamente, se asume que el movimiento ocurre en un régimen dominado por fuerzas viscosas, donde la velocidad terminal se alcanza cuando la fuerza de arrastre equilibra el peso efectivo de la esfera (considerando el empuje). Asimismo, se supone que el flujo alrededor de la esfera es aproximadamente laminar, que la viscosidad del aceite se mantiene constante durante cada experimento, y que los efectos de borde del recipiente no alteran significativamente el comportamiento dinámico del sistema. Bajo estas condiciones, se espera observar un comportamiento lineal en la relación altura-tiempo durante el régimen terminal, lo que permitiría estimar la velocidad constante del sistema y contrastarla con las predicciones del modelo físico planteado. 

\indent Procedemos a exponer los \secref{sec:RESULTADOS} obtenidos. Si nos guiamos por la \secref{sec:METODOS}, se hicieron tres réplicas para cada tipo de canica\footnote{Recordemos que hay tres tipos o categorías de canicas (esferas): la canica negra, la canica blanca, y la canica de metal. Esta distinción de categorización se basa en la masa de cada una de la canicas. La masas están catalogadas en \secref{sec:DATOS}.}, donde usando \Trackercode,  se pudo obtener el tiempo y la altura\footnote{Notemos que el punto de referencia fue en el eje negativo, y es por eso que las alturas están negativas.}. Para aprovechar las tres replicas se implemento una regresión lineal con efectos fijos por grupo \citet{christensen2020analysis}; se propuso este modelo debido a que se considero que los coeficientes de las velocidades terminales para cada canica es constante y al promediar las tres replicas por Mínimos Cuadrados Ordinarios (OLS) usualmente debería dar una mejor estimación para los coeficientes.\newline

\indent Los resultados por la regresión lineal con efectos fijos por grupos se pueden observar en \cref{tab:II}. En el cuadrante de los coeficientes, la primera columna a mano izquierda, tiene como primera indicación la variable independiente \textit{Tiempo} con las tres replicas agrupadas (la velocidad terminal promediada $\hat{v}_t$). Los dos renglones que le proceden son la diferencia entre $\hat{v}_t$ y la velocidad terminal obtenida en la replica $\hat{v}_{t,2}, \hat{v}_{t,3}$\footnote{Un detalle a mencionar es que usar un modelo con efectos fijos requiere usar una replica como la replica referencial, y es por eso que la replica numero uno no aparece en la tabla.}. Lo importante a notar es que todas las variables son significativas con $(\text{P-value}<0.01)$, y además la diferencia entre $\hat{v}_t$ y las otras replicas son indistinguibles por lo que podemos sustentar que los coeficientes $\hat{v}_t$ se aproxima adecuadamente a su valor teórico. Además, debajo de los coeficientes\footnote{También se quisiera aclarar un detallo sobre el numero de observaciones. Debido a que se agruparon las replican en un solo modelo, es que consecuentemente también se acumularon todas las observaciones para estimar $\hat{v}_t$.  Otro detalle es que aquí las observaciones indican en numero de fotogramas que se logro capturar por el video.} podemos observar que tanto el coeficiente de determinación y la prueba de significancia global ($F$-statistic) muestran un buen ajuste para una relación lineal entre el tiempo y la altura. Aunque se debería de notar que estos resultados mencionado anteriormente no tienen significancia para validar la certeza y consistencia de los coeficientes obtenidos, es por ello que a continuación se va a explicar el diagnostico de los residuos. \newline

\indent Las tres tablas que se muestran debajo de los coeficientes en \cref{tab:II} indican si cada modelo para la canicas cumple los supuestos de los residuos\footnote{Debido a que usamos una regresión con efectos fijos, se puede omitir probar la suposición de Media Zero, osea que el valor esperado de los residuos es zero.}, siendo estos normalidad, Homocedasticidad, y independencia. Como se puede observar, todos los modelos agrupados para la canica blanca, negra y de metal no cumplen ninguno de los supuestos. Aunque notemos que como la variable independiente es el tiempo, era muy factible que los residuos iban a experimentar un nivel de dependencia \citet{christensen2020analysis}. Entonces, los supuestos de normalidad y homocedasticidad son los supuestos que esta discusión se va a enfocar en exponer y documentar las consecuencia que estos supuestos hayan fallado en la metodología propuesta.\newline

\indent Para observar el diagnóstico de residuos, se presentan en \secref{sec:Graficos} para cada canica cuatro diferentes gráficas sobre los residuos. Para esto observemos los residuos para la canica de metal (\cref{fig:3}). Algo importante a observar es que gráficamente podemos deducir que en general el modelo de la canica de metal tiene un comportamiento lineal, con un puntos de influencia por parte de la observación numero 7. Ademas podemos ver que la curva de suavizado (LOESS)\footnote{Recordatorio: la curva de suavizado (LOESS) en gráficos de residuos es un ajuste local que resume la tendencia sistemática de los residuos. En el lenguaje de programación $R$ se presenta como una linea roja en gráficos de residuos (\cite{christensen2020analysis}).} en general tiende a seguir una trayectoria lineal lo que nos afirma que la aceleración experimentada por la canica de metal era despreciable. Dicho lo anterior, es muy factible que al remover los puntos de influencia, y limpiar los datos, la regresión para la canica de metal sí satisfacería los supuestos ante los residuos. Pero, como se discutirá a continuación, esto ya no se llevó a cabo debido a lo que se observó en las otras gráficas de diagnósticos.\newline

\indent Las gráficas \cref{fig:4} y \cref{fig:5} son los diagnósticos para la canica negra y la canica blanca respectivamente. Al comparar estas gráficas con \cref{fig:3}, es bastante evidente que las canicas blancas y negras no tienen un diagnostico de residuos favorable como la tuvo la canica de metal. Queremos nuevamente recordar al lector que estas tres gráficas son el modelo agrupado con las tres replicas, y es por ello que uno puede observar diferentes trayectorias o pendientes en la misma gráfica. Estas diferentes trayectorias se pueden visualizar en \cref{fig:4} para la imagen de la fila superior izquierda, donde se presenta que entre las replicas (de la misma canica) indican diferentes pendiente para la LOESS. En especifico la gráfica de la fila inferior a la izquierda, indica que hay una curvatura notable para el LOESS. Esto nos dice que la canica negra y la blanca, en el intervalo de tiempo considerado, experimentaron aceleración no despreciable, y por ello los residuos no condicen con los valores ajustados.\newline

Recopilando respecto los diagnósticos de residuos, no podemos estudiar la \secref{sec:HIPOTESIS} de este experimento debido a que dos de tres observaciones tienen defectos de estudio. Por ello no podemos garantizar si los coeficientes de la velocidad terminal son confiables para trabajar; nuevamente, con la excepción de la canica de metal la que requiere el uso de limpieza de datos para determinar la certeza del coeficiente. Algo que podemos afirmar, en parte por lo que se presentó en \secref{sec:MacroTeorico}, es que la metodología propuesta es viable si se usan masas lo suficientemente pesadas para que se aproximen adecuadamente a su velocidad terminal, como se observo en la canica de metal. Por otro lado, con la canica negra y blanca, al tener una masa muy liviana, no se llegó a obtener coeficientes confiables.

En resumen, las canicas más livianas no desaceleran lo suficientemente rápido para que sea válido estimar la velocidad terminal mediante una regresión lineal de los datos de posición y tiempo. En consecuencia, se recomienda el uso de masas más pesadas, aunque se reconoce la dificultad de variar la masa sin variar el volumen de las esferas, siendo esta la principal debilidad de la metodología propuesta. Una solución alternativa sería utilizar solo una parte de los datos, en donde las canicas hayan desacelerado lo suficiente, como se propone en \citet{gluck2003air}; aunque como se expone en los antecedentes esto introduciría sesgo. Una tercera alternativa sería realizar una regresión sobre los datos de posición tiempo de la solución analítica expresada en el marco teórico: \eqref{eq:x-t}

















%Como se ve en la \cref{tab:II},
%\cref{hyp:A}
%\hyperref[obj:A]{Objetivo A}

%for apa7 reference use: \parencite{smith2019study}  (author, date)
%to use in text citation use: \textcite{garcia2020efecto}. author (date)

%To reference an objective do~\hyperref[obj:A]{Objetivo A}
%To reference a equation do~\autoref{eq:equationLabel} 

%To reference a table use ~\autoref{tab:TableNO} this prints the label and a hyperlink to that table


%To reference a figure use ~\autoref{fig:figNo} this prints the figure number and the hyperlink
\section{Conclusiones}
\label{sec:CONCLUSIONES}

\begin{itemize}
    \item Con los datos obtenidos, no hay confiabilidad para estudiar la linealidad entre la velocidad terminal y la masa del objecto en una caída libre usando una viscosidad de aceite comercial.
    \item La metodología propuesta solo permite la confiabilidad del coeficiente de velocidad terminal para objectos pesado como la canica de metal ($16.311$ $g$). Mientras que objectos livianos, como la canica negra ($5.324$ $g$) y blanca ($5.316$ $g$) surgió en los diagnósticos de residuos  desconfianza ante los coeficientes del modelo. 
    \item Se sugiere para futuras replicas de esta metodología usar masas similares a la canica de metal para garantizar confiabilidad de los resultados. Ademas sugerimos investigar la masa mínima para garantizar confiabilidad en una metodología similar a esta para evitar situaciones similar a la que sucedió en este experimento. 
\end{itemize}




\footnotesize
\printbibliography
\normalsize



\section*{Anexo}
\label{sec:Anexo}

%param: Title, Code, Hyperlink
\begin{FigureAnexo}{Tabla de Nutrientes del Aceite}{}{figA:tabla aceite}
\includegraphics[width=0.5\linewidth]{images/practica_1/aceite.png}
\end{FigureAnexo}



\end{document}
 